\documentclass[12pt]{article}
\usepackage{fullpage,amsmath,amssymb,amsthm,hyperref,amsfonts}
\newtheorem{theorem}{Theorem}
\newtheorem{lemma}[theorem]{Lemma}
\newtheorem{proposition}[theorem]{Proposition}
\theoremstyle{definition}
\newtheorem{definition}[theorem]{Definition}
\newtheorem{remark}[theorem]{Remark}

\title{Lagrangian Smoothings of Polyhedral Lagrangian Surfaces \\ with $4$-Valent Vertices}
\author{}
\date{}

\begin{document}
\maketitle

\section{Statement and Result}

Throughout, $(\mathbb{R}^4,\omega_0)$ denotes the standard symplectic vector space, with coordinates $(x_1,y_1,x_2,y_2)$ and symplectic form
\begin{equation}\label{eq:omega}
\omega_0 \;=\; dx_1\wedge dy_1 + dx_2\wedge dy_2.
\end{equation}
A linear subspace $\Lambda\subset\mathbb{R}^4$ of dimension $2$ is \emph{Lagrangian} if $\omega_0|_\Lambda=0$. A smooth $2$-dimensional submanifold $L\subset\mathbb{R}^4$ is Lagrangian if every tangent plane $T_pL$ is a Lagrangian subspace.

\begin{definition}[Polyhedral Lagrangian surface]\label{def:polyhedral}
A \emph{polyhedral Lagrangian surface} $K\subset(\mathbb{R}^4,\omega_0)$ is a finite polyhedral $2$-complex
\[
K \;=\; \bigcup_{F\in\mathcal{F}} F
\]
in which each face $F\in\mathcal{F}$ is a Lagrangian polygon (a flat $2$-polygon contained in some affine Lagrangian plane), such that the underlying topological space $|K|$ is a topological submanifold of $\mathbb{R}^4$.
\end{definition}

\begin{definition}[Lagrangian smoothing]\label{def:smoothing}
A \emph{Lagrangian smoothing} of a polyhedral Lagrangian surface $K$ is a family $(K_t)_{t\in[0,1]}$ of subsets of $\mathbb{R}^4$ such that:
\begin{enumerate}
\item For each $t\in(0,1]$, $K_t$ is a smooth Lagrangian submanifold of $(\mathbb{R}^4,\omega_0)$.
\item The family $(K_t)_{t\in(0,1]}$ is a Hamiltonian isotopy: there exists a smooth time-dependent Hamiltonian $H:\mathbb{R}^4\times(0,1]\to\mathbb{R}$, compactly supported, whose Hamiltonian flow $\Phi_t^H$ satisfies $\Phi_t^H(K_1)=K_t$ for $t\in(0,1]$.
\item The map $t\mapsto K_t$ extends continuously to $[0,1]$ with respect to the Hausdorff topology on compact subsets, and $K_0=K$.
\item The full family $(K_t)_{t\in[0,1]}$ is a topological isotopy: there exists a continuous map $\Psi:|K|\times[0,1]\to\mathbb{R}^4$ with $\Psi(\cdot,t)$ a homeomorphism onto $K_t$ and $\Psi(\cdot,0)=\mathrm{id}_{|K|}$.
\end{enumerate}
\end{definition}

The vertex set $V(K)$ of $K$ is the $0$-skeleton of the polyhedral structure. We say that $K$ is \emph{$4$-valent} if every vertex $p\in V(K)$ is contained in exactly $4$ faces.

\begin{theorem}[Main result]\label{thm:main}
Let $K\subset(\mathbb{R}^4,\omega_0)$ be a polyhedral Lagrangian surface in which exactly $4$ faces meet at every vertex. Then $K$ admits a Lagrangian smoothing in the sense of Definition~\ref{def:smoothing}.
\end{theorem}

The answer to Abouzaid's question is therefore \textbf{yes}.

\medskip

The proof proceeds by constructing $(K_t)$ explicitly. The skeleton of the construction is:
\begin{enumerate}
\item Identify the local Lagrangian cone $C_p$ at each vertex $p$. Show that the link $L_p = K\cap S^3_\varepsilon(p)$ is a piecewise-Legendrian topological unknot in the standard contact $(S^3_\varepsilon,\xi_{\mathrm{std}})$ with classical invariants $\mathrm{tb}(L_p)=-1$ and $r(L_p)=0$ (Lemma~\ref{lem:link} and Lemma~\ref{lem:invariants}).
\item Conclude that $L_p$ bounds a smooth Lagrangian disk in the symplectic $4$-ball, by the Eliashberg--Fraser classification of Legendrian unknots (Lemma~\ref{lem:eliashberg-fraser}).
\item Replace the cone $C_p$ inside a small ball $B_p$ by a smooth Lagrangian disk filling $L_p$, using a compactly supported Hamiltonian isotopy that extends continuously to $t=0$ (Proposition~\ref{prop:vertex}).
\item Smooth the polyhedral structure along edges by a one-parameter family of Hamiltonian rotations (Lemma~\ref{lem:edge}).
\item Compose the disjoint local isotopies into a global Hamiltonian isotopy $\Phi_t$ of $\mathbb{R}^4$ and set $K_t=\Phi_t(K)$ (proof of Theorem~\ref{thm:main}).
\end{enumerate}

\section{The Local Cone at a $4$-Valent Vertex}

\subsection{The tangent cone}

Fix a vertex $p\in V(K)$. The four Lagrangian faces meeting at $p$ have well-defined tangent half-planes at $p$; we denote them $H_1,H_2,H_3,H_4$, each a closed Lagrangian half-plane in $T_p\mathbb{R}^4\cong\mathbb{R}^4$. Let $\Lambda_i\subset T_p\mathbb{R}^4$ be the (full) Lagrangian plane containing $H_i$.

\begin{definition}[Tangent cone]\label{def:cone}
The \emph{tangent cone} of $K$ at $p$ is
\begin{equation}\label{eq:cone}
C_p \;=\; H_1\cup H_2\cup H_3\cup H_4 \;\subset\; T_p\mathbb{R}^4.
\end{equation}
\end{definition}

Since $|K|$ is a topological $2$-manifold near $p$, the union $C_p\setminus\{0\}$ is a topological $2$-manifold; equivalently, the link
\[
L_p^{\mathrm{cone}} \;=\; C_p \cap S^3_1(0)
\]
is a topological $1$-manifold in the unit sphere $S^3_1(0)\subset T_p\mathbb{R}^4$. As $L_p^{\mathrm{cone}}$ is the union of four arcs $A_i = H_i\cap S^3_1(0)$, each homeomorphic to an interval, and these arcs glue into a $1$-manifold without boundary, the topology forces them to assemble into a single closed loop, i.e.\ a topological circle.

We record this combinatorial constraint formally.

\begin{lemma}[Cyclic gluing]\label{lem:cyclic}
There exists a cyclic ordering of the half-planes, which we may take to be $H_1,H_2,H_3,H_4$, such that the boundary rays satisfy
\begin{equation}\label{eq:cyclic}
\partial H_i \cap \partial H_{i+1} \neq \{0\} \qquad (i\in\mathbb{Z}/4\mathbb{Z}),
\end{equation}
and the link $L_p^{\mathrm{cone}}$ is the topological circle obtained by joining the arcs $A_1,A_2,A_3,A_4$ in cyclic order.
\end{lemma}

\begin{proof}
Each half-plane $H_i$ has two boundary rays through the origin (its two sides as a half-plane are not present here; rather, $\partial H_i$ is a single line; but the half-plane $H_i$ has a one-dimensional boundary line $\partial H_i$ together with the apex $0$). In the link $L_p^{\mathrm{cone}}=C_p\cap S^3_1(0)$, each arc $A_i$ has two endpoints, namely the two points $\partial H_i\cap S^3_1(0)$. Since $L_p^{\mathrm{cone}}$ is a topological $1$-manifold without boundary, every endpoint of every arc must coincide with the endpoint of exactly one other arc. The endpoints thus pair up; the pairing partitions the eight endpoints into four pairs and induces a graph on the index set $\{1,2,3,4\}$ in which every vertex has degree $2$. Such a graph is either a $4$-cycle or two disjoint $2$-cycles. The latter would force the union $A_i\cup A_j$ to be a closed loop disjoint from $A_k\cup A_\ell$, contradicting connectedness of $L_p^{\mathrm{cone}}$ (which is required since $|K|$ is a $2$-manifold near $p$, so $C_p\setminus\{0\}$ is connected). Hence the graph is a single $4$-cycle, so the arcs assemble cyclically.
\end{proof}

\subsection{The model product cone}

\begin{definition}[Standard product cone]\label{def:product-cone}
The \emph{standard product cone} $C_0\subset\mathbb{R}^4$ is defined by
\begin{equation}\label{eq:product-cone}
C_0 \;=\; \Gamma\times\Gamma, \qquad \Gamma \;=\; \{(x,0):x\geq 0\}\cup\{(0,y):y\geq 0\}\subset\mathbb{R}^2,
\end{equation}
where the product is taken under the identification $\mathbb{R}^4\cong\mathbb{R}^2\times\mathbb{R}^2$ with coordinates $((x_1,y_1),(x_2,y_2))$.
\end{definition}

The set $C_0$ is the union of four Lagrangian half-planes
\[
H_1^{(0)}=\{x_1\geq 0,x_2\geq 0,y_1=y_2=0\},\quad H_2^{(0)}=\{x_1\geq 0,y_2\geq 0,y_1=x_2=0\},
\]
\[
H_3^{(0)}=\{y_1\geq 0,y_2\geq 0,x_1=x_2=0\},\quad H_4^{(0)}=\{y_1\geq 0,x_2\geq 0,x_1=y_2=0\}.
\]
Each $H_i^{(0)}$ is contained in a coordinate Lagrangian plane (e.g.\ $H_1^{(0)}\subset\{y_1=y_2=0\}$, etc.), and they meet cyclically along the four positive coordinate rays.

\begin{lemma}[Linear straightening]\label{lem:straighten}
There exists a linear symplectomorphism $A_p\in\mathrm{Sp}(4,\mathbb{R})$ that maps the tangent cone $C_p$ at $p$ to the standard product cone $C_0$, and a connecting smooth path of linear symplectomorphisms $A_p^{(s)}$, $s\in[0,1]$, with $A_p^{(0)}=A_p$ and $A_p^{(1)}\cdot C_0 = C_p$.
\end{lemma}

\begin{proof}
We claim that the configuration space
\[
\mathcal{V} \;=\; \bigl\{(\Lambda_1,\Lambda_2,\Lambda_3,\Lambda_4)\in(\mathrm{LGr}(4))^4 \mid \Lambda_i\cap\Lambda_{i+1}\neq\{0\},\;\Lambda_i\cap\Lambda_{i+2}=\{0\}\bigr\}
\]
of cyclically ordered $4$-tuples of Lagrangian planes meeting in lines and pairwise transverse otherwise is connected, and that it acts transitively under the symplectic group $\mathrm{Sp}(4,\mathbb{R})$.

\smallskip\noindent\emph{Transitivity.} Given $(\Lambda_1,\Lambda_2,\Lambda_3,\Lambda_4)\in\mathcal{V}$, choose unit vectors $v_i\in\Lambda_i\cap\Lambda_{i+1}$ for $i=1,2,3,4$ (indices mod $4$); these four vectors generate the cone since $\Lambda_i=\mathrm{span}(v_{i-1},v_i)$ (as $v_{i-1},v_i\in\Lambda_i$ and they span a $2$-plane because $\Lambda_{i-1}\cap\Lambda_{i+1}=\{0\}$ forces $v_{i-1}\notin\Lambda_{i+1}$, so $v_{i-1}\not\parallel v_i$).

The Lagrangian condition $\omega_0|_{\Lambda_i}=0$ gives
\begin{equation}\label{eq:omega-vij}
\omega_0(v_{i-1},v_i)=0 \quad (i=1,2,3,4).
\end{equation}
The transversality $\Lambda_i\cap\Lambda_{i+2}=\{0\}$ implies the spans $\mathrm{span}(v_{i-1},v_i)$ and $\mathrm{span}(v_{i+1},v_{i+2})$ are transverse, hence $v_1,v_2,v_3,v_4$ form a basis of $\mathbb{R}^4$. The remaining symplectic pairings $\omega_0(v_i,v_{i+2})$ for $i=1,2$ are free: by rescaling $v_i$ we can arrange $|\omega_0(v_1,v_3)|=|\omega_0(v_2,v_4)|=1$. After possibly replacing $v_i$ by $-v_i$ and reversing cyclic orientation, the standard model has
\[
v_1=\partial_{x_1},\quad v_2=\partial_{x_2},\quad v_3=\partial_{y_1},\quad v_4=\partial_{y_2},
\]
which satisfies $\omega_0(v_1,v_3)=1$, $\omega_0(v_2,v_4)=1$, $\omega_0(v_i,v_{i\pm 1})=0$. The unique linear map sending the original basis $(v_1,v_2,v_3,v_4)$ to the model basis $(\partial_{x_1},\partial_{x_2},\partial_{y_1},\partial_{y_2})$ then preserves $\omega_0$, hence lies in $\mathrm{Sp}(4,\mathbb{R})$.

\smallskip\noindent\emph{Connectivity.} The map $A_p^{(s)} = A_p\cdot\exp(s\cdot 0)= A_p$ deforms within $\mathrm{Sp}(4,\mathbb{R})$, which is connected (it is homotopy-equivalent to $U(2)$). Choosing any continuous path in $\mathrm{Sp}(4,\mathbb{R})$ from $\mathrm{Id}$ to $A_p^{-1}$ yields the required path.

The orientation of the cyclic ordering may be reversed by the symplectomorphism $(x_1,y_1,x_2,y_2)\mapsto(y_1,x_1,y_2,x_2)$, which preserves $\omega_0$ up to sign and can be composed with reflection to remain in $\mathrm{Sp}(4,\mathbb{R})$.
\end{proof}

\subsection{Standard contact structure on the link sphere}

For $\varepsilon>0$ small, denote the sphere $S^3_\varepsilon(p)=\{q\in\mathbb{R}^4:|q-p|=\varepsilon\}$. Translating $p$ to the origin and identifying $\mathbb{R}^4=\mathbb{C}^2$ via $z_1=x_1+iy_1$, $z_2=x_2+iy_2$, the standard contact form is
\begin{equation}\label{eq:alpha}
\alpha_{\mathrm{std}} \;=\; \tfrac{1}{2}\bigl(x_1\,dy_1-y_1\,dx_1+x_2\,dy_2-y_2\,dx_2\bigr)\Big|_{S^3_\varepsilon},
\end{equation}
with contact distribution $\xi_{\mathrm{std}}=\ker\alpha_{\mathrm{std}}$. The form $\alpha_{\mathrm{std}}$ is the restriction of the Liouville form $\lambda_0=\tfrac{1}{2}\sum(x_i\,dy_i-y_i\,dx_i)$, which satisfies $d\lambda_0=\omega_0$.

\begin{lemma}[Lagrangian cones meet $S^3_\varepsilon$ Legendrian-ly]\label{lem:legendrian}
Let $\Lambda\subset\mathbb{R}^4$ be a Lagrangian half-plane through the origin and $0<\varepsilon<\infty$. The intersection $\Lambda\cap S^3_\varepsilon(0)$ is a smooth arc, every tangent line of which lies in the contact distribution $\xi_{\mathrm{std}}$.
\end{lemma}

\begin{proof}
Parametrize $\Lambda\cap S^3_\varepsilon$ by arc length $\gamma(s)$. The tangent vector $\gamma'(s)$ lies in $T_{\gamma(s)}\Lambda=\Lambda$ (since $\Lambda$ is a linear subspace) and is perpendicular to the radial vector $\gamma(s)$ (since $\gamma$ lies on $S^3_\varepsilon$). The Liouville form satisfies
\[
\lambda_0(\gamma(s))(v) \;=\; \tfrac{1}{2}\omega_0(\gamma(s),v) \qquad\text{for any } v\in T_{\gamma(s)}\mathbb{R}^4,
\]
because $\lambda_0(q)(v)=\tfrac{1}{2}\sum(x_i\,dy_i(v)-y_i\,dx_i(v))=\tfrac{1}{2}\omega_0(q,v)$ where on the right we view $q$ as a tangent vector at $q$. Hence
\[
\alpha_{\mathrm{std}}(\gamma'(s)) \;=\; \tfrac{1}{2}\omega_0(\gamma(s),\gamma'(s)) \;=\; 0,
\]
since both $\gamma(s)$ (radial direction in $\Lambda$) and $\gamma'(s)$ lie in $\Lambda$ and $\omega_0|_\Lambda=0$. Thus $\gamma'(s)\in\xi_{\mathrm{std}}$.
\end{proof}

\begin{lemma}[Link is a piecewise-Legendrian unknot]\label{lem:link}
For every sufficiently small $\varepsilon>0$, the intersection
\begin{equation}\label{eq:link}
L_p \;=\; K\cap S^3_\varepsilon(p)
\end{equation}
is a piecewise-Legendrian topological circle in $(S^3_\varepsilon,\xi_{\mathrm{std}})$, which is unknotted (i.e.\ bounds an embedded topological disk).
\end{lemma}

\begin{proof}
By Definition~\ref{def:polyhedral}, $|K|$ is a topological $2$-manifold near $p$, so for $\varepsilon$ small enough that no other vertex or non-radial part of $K$ enters the closed ball $\overline{B}_\varepsilon(p)$, the intersection $K\cap\overline{B}_\varepsilon(p)$ is a topological cone over $L_p=K\cap S^3_\varepsilon(p)$. Since $|K|$ is a $2$-manifold, the cone is a $2$-disk, hence $L_p$ is a topological $1$-sphere bounding the disk $K\cap\overline{B}_\varepsilon(p)$ in $\overline{B}_\varepsilon(p)$.

For $\varepsilon$ small, $K\cap\overline{B}_\varepsilon(p)$ approximates the linear cone $C_p$ in $C^0$, and the four arcs $A_i=H_i\cap S^3_\varepsilon$ are smooth (each is a Legendrian arc on the geodesic of $S^3_\varepsilon$ in the Lagrangian plane $\Lambda_i$). By Lemma~\ref{lem:cyclic}, these arcs concatenate cyclically into the topological $1$-sphere $L_p^{\mathrm{cone}} = C_p\cap S^3_\varepsilon$. The actual link $L_p=K\cap S^3_\varepsilon$ differs from $L_p^{\mathrm{cone}}$ by a small Hausdorff-distance perturbation that respects the polyhedral structure: each arc of $L_p$ is the intersection of one of the four Lagrangian faces $F_i$ with $S^3_\varepsilon$, and by Lemma~\ref{lem:legendrian} this arc is Legendrian.

Unknottedness: $L_p$ bounds the disk $K\cap\overline{B}_\varepsilon(p)$ in $\overline{B}_\varepsilon(p)$, hence in $S^3_\varepsilon$ as the boundary of a punctured tubular neighborhood; explicitly, push the disk radially outward to a disk in $S^3_\varepsilon$ with the same boundary $L_p$. Thus $L_p$ is the unknot.
\end{proof}

\subsection{Classical invariants of the link}

For a smooth Legendrian knot $L\subset(S^3,\xi_{\mathrm{std}})$, the \emph{Thurston--Bennequin number} $\mathrm{tb}(L)$ is the framing of $L$ given by $\xi_{\mathrm{std}}$ minus the Seifert framing. The \emph{rotation number} $r(L)$ is the winding number of $TL$ in $\xi_{\mathrm{std}}|_L$ relative to a trivialization extending over a Seifert surface. Both are integers, well-defined up to homotopy of the knot through Legendrians.

For a piecewise-Legendrian knot (with finitely many corners) we use the following extension. Each arc $A_i$ is Legendrian smooth, and at each corner the two adjacent Legendrian directions span a $2$-plane in $\xi_{\mathrm{std}}$ along which the corner is rounded by an arbitrarily small Legendrian arc; the resulting smooth Legendrian is well-defined up to Legendrian isotopy after the rounding choice, and $\mathrm{tb}$, $r$ are independent of the choice.

\begin{lemma}[Computation for the product cone]\label{lem:invariants}
Let $L_0=C_0\cap S^3_\varepsilon$ be the link of the standard product cone~\eqref{eq:product-cone}. After Legendrian rounding at the four corners, $L_0$ is Legendrian-isotopic to the standard Legendrian unknot $U_0\subset(S^3,\xi_{\mathrm{std}})$ with
\begin{equation}\label{eq:tb-r}
\mathrm{tb}(L_0)=-1,\qquad r(L_0)=0.
\end{equation}
\end{lemma}

\begin{proof}
Decompose $\mathbb{R}^4=\mathbb{R}^2\oplus\mathbb{R}^2$ with $\omega_0=\omega_1\oplus\omega_2$, where $\omega_i=dx_i\wedge dy_i$. The cone $C_0=\Gamma\times\Gamma$ is a product of two V-shaped Lagrangian curves. We first smooth $\Gamma$ in $\mathbb{R}^2$ to a Lagrangian curve $\gamma$.

Since any smooth curve in $\mathbb{R}^2$ is automatically Lagrangian (every $2$-form on a $1$-manifold vanishes), we choose $\gamma\subset\mathbb{R}^2$ to be the smooth quarter-circle $\gamma=\{(r\cos\theta,r\sin\theta):\theta\in[0,\pi/2]\}$ for some fixed $r>0$. Then $\widetilde{C}_0=\gamma\times\gamma\subset\mathbb{R}^4$ is a smooth Lagrangian disk of dimension $2$, since the product of two Lagrangians in two symplectic factors is Lagrangian in the direct sum.

The boundary $\partial\widetilde{C}_0$ consists of four Legendrian arcs (smoothed corners) and lies on the $3$-sphere of radius $r\sqrt{2}$ in $\mathbb{R}^4$. After rounding the four corners (where $\gamma\times\{(r,0)\}$ meets $\gamma\times\{(0,r)\}$, etc.) by a Legendrian isotopy, $\partial\widetilde{C}_0$ becomes a smooth Legendrian unknot in $S^3$.

We now apply the Cieliebak--Eliashberg formula for the classical invariants of the boundary of an exact Lagrangian filling: if $D\subset(B^4,\omega_0)$ is a smooth exact Lagrangian disk with Legendrian boundary $L\subset(S^3,\xi_{\mathrm{std}})$, then
\begin{equation}\label{eq:CE}
\mathrm{tb}(L) \;=\; -\chi(D),\qquad r(L) \;=\; \mu(D)/2,
\end{equation}
where $\chi(D)$ is the Euler characteristic and $\mu(D)$ the Maslov index of $D$ (relative to a Lagrangian framing on the boundary).

\smallskip\noindent\emph{Hypothesis verification for~\eqref{eq:CE}.} The Lagrangian $\widetilde{C}_0$ is exact: the Liouville form $\lambda_0=\tfrac{1}{2}\sum(x_i\,dy_i-y_i\,dx_i)$ pulled back to $\widetilde{C}_0=\gamma\times\gamma$ splits as the sum of the pullbacks of $\lambda_i=\tfrac{1}{2}(x_i\,dy_i-y_i\,dx_i)$ to the two $\gamma$-factors. Each pullback to a smooth curve in $\mathbb{R}^2$ is a smooth $1$-form, hence is automatically exact on the disk-shaped factor (which is contractible). So $\lambda_0|_{\widetilde{C}_0}$ is exact.

\smallskip\noindent\emph{Computation of $\chi$.} As $\widetilde{C}_0=\gamma\times\gamma$ is a topological square, $\chi(\widetilde{C}_0)=1$, so $\mathrm{tb}(L_0)=-1$.

\smallskip\noindent\emph{Computation of $\mu$.} The Maslov index $\mu(\widetilde{C}_0)$ is the signed intersection number of the Gauss map $\widetilde{C}_0\to\mathrm{LGr}(4)$ with the Maslov cycle. Since $\widetilde{C}_0=\gamma\times\gamma$ is a product, its Gauss map factors as a product of the Gauss maps $\gamma\to\mathrm{LGr}(2)\cong S^1$. Each factor $\gamma$ is a quarter-arc of the circle, so its Gauss map covers a quarter of $S^1$; the product Gauss map is thus a small disk in $\mathrm{LGr}(4)$, contractible to a point. Hence $\mu(\widetilde{C}_0)=0$, giving $r(L_0)=0$.

\smallskip\noindent\emph{Identification with $U_0$.} The standard Legendrian unknot $U_0$ in $(S^3,\xi_{\mathrm{std}})$ is characterized by $\mathrm{tb}=-1,r=0$, and bounds the standard Lagrangian disk $D_0$ given by the real $2$-plane in the symplectic $4$-ball after suitable rotation. By the Eliashberg--Fraser theorem (Lemma~\ref{lem:eliashberg-fraser} below), Legendrian unknots with $\mathrm{tb}=-1,r=0$ are unique up to Legendrian isotopy. Hence $L_0$ is Legendrian-isotopic to $U_0$.
\end{proof}

\begin{lemma}[Constancy across $\mathcal{V}$]\label{lem:constancy}
Let $\mathcal{V}^-\subset\mathcal{V}$ be the connected component of $\mathcal{V}$ (Lemma~\ref{lem:straighten}) corresponding to the cyclic ordering matching that of the standard product cone. For every $(\Lambda_1,\dots,\Lambda_4)\in\mathcal{V}^-$, the link $L_p$ of the corresponding cone $C_p=H_1\cup\dots\cup H_4$ in $S^3_\varepsilon$ has, after Legendrian rounding at the corners,
\begin{equation}\label{eq:tb-r-general}
\mathrm{tb}(L_p)=-1,\qquad r(L_p)=0.
\end{equation}
\end{lemma}

\begin{proof}
By Lemma~\ref{lem:straighten}, every $(\Lambda_1,\dots,\Lambda_4)\in\mathcal{V}^-$ is the image of $(\Lambda_1^{(0)},\dots,\Lambda_4^{(0)})$ under some $A\in\mathrm{Sp}(4,\mathbb{R})$. The action of $A$ on $S^3_\varepsilon$ is a contactomorphism after rescaling, hence preserves $\mathrm{tb}$ and $r$ of any Legendrian. Lemma~\ref{lem:invariants} gives $\mathrm{tb}=-1,r=0$ in the standard model; transport gives the same in general.
\end{proof}

\subsection{Eliashberg--Fraser classification}

\begin{lemma}[Eliashberg--Fraser]\label{lem:eliashberg-fraser}
Any Legendrian unknot $L\subset(S^3,\xi_{\mathrm{std}})$ with $\mathrm{tb}(L)=-1$ and $r(L)=0$ is Legendrian-isotopic to the standard Legendrian unknot $U_0$.
\end{lemma}

\noindent\emph{Reference.} Eliashberg--Fraser, \emph{Classification of topologically trivial Legendrian knots}, CRM Proc.\ Lecture Notes 15 (1998), 17--51, Theorem~1.2.

\noindent\emph{Hypothesis verification.} The hypotheses of the theorem are: (i) $L$ is topologically the unknot; (ii) $L$ is Legendrian for $\xi_{\mathrm{std}}$; (iii) $\mathrm{tb}(L)=-1,r(L)=0$. By Lemma~\ref{lem:link}, $L_p$ is unknotted. By Lemma~\ref{lem:legendrian} (after Legendrian rounding), $L_p$ is Legendrian. By Lemma~\ref{lem:constancy}, the invariants are as required. So the theorem applies.

\begin{lemma}[Standard Lagrangian disk]\label{lem:disk}
The standard Legendrian unknot $U_0$ in $(S^3,\xi_{\mathrm{std}})$ bounds a smooth embedded Lagrangian disk $D_0$ in the symplectic $4$-ball $(B^4,\omega_0)$.
\end{lemma}

\begin{proof}
Take $D_0=\{(x_1+iy_1,x_2+iy_2)\in B^4:y_1=y_2=0\}=\mathbb{R}^2\cap B^4$. Then $\omega_0|_{D_0}=dx_1\wedge dy_1+dx_2\wedge dy_2$ vanishes (no $dy_i$ pulled back since $y_i=0$ on $D_0$). Its boundary $\partial D_0=\{y_1=y_2=0\}\cap S^3$ is the standard Legendrian unknot, after a small contact-isotopy if needed; equivalently, $\partial D_0$ is the maximal Legendrian unknot in the front projection (a circle of front-radius $1$). Cf.\ Etnyre, \emph{Legendrian and transversal knots}, Hdb.\ Knot Theory (2005), §2.7.
\end{proof}

\section{Local Hamiltonian Smoothing at a Vertex}

\begin{proposition}[Vertex smoothing]\label{prop:vertex}
Let $p\in V(K)$ be a $4$-valent vertex with cone $C_p$. There exists $\varepsilon>0$ and a one-parameter family $(\Phi_t^{(p)})_{t\in[0,1]}$ of homeomorphisms of $\mathbb{R}^4$ supported in $\overline{B}_\varepsilon(p)$ such that:
\begin{enumerate}
\item For each $t\in(0,1]$, $\Phi_t^{(p)}$ is the time-$1$ map of a smooth Hamiltonian $H_t:\mathbb{R}^4\to\mathbb{R}$ with $\mathrm{supp}(H_t)\subset\overline{B}_\varepsilon(p)$, depending smoothly on $t$.
\item For $t=1$, $\Phi_1^{(p)}=\mathrm{id}$.
\item For each $t\in(0,1]$, $\Phi_t^{(p)}(K\cap B_\varepsilon(p))$ is smooth Lagrangian.
\item The family $t\mapsto\Phi_t^{(p)}$ is continuous on $[0,1]$ in the $C^0$ (Hausdorff) topology, with $\Phi_0^{(p)}=\mathrm{id}$.
\end{enumerate}
\end{proposition}

\begin{proof}
Choose $\varepsilon>0$ so small that $K\cap\overline{B}_\varepsilon(p)$ coincides with $C_p\cap\overline{B}_\varepsilon(p)$ (possible because $K$ is locally polyhedral with cone $C_p$ at $p$ — every face is flat near $p$). By Lemma~\ref{lem:straighten}, after composing with a linear symplectomorphism we may assume $C_p=C_0$ (the standard product cone) and $p=0$. We construct the family on this normal form; conjugating back by the linear symplectomorphism gives the result for general $C_p$.

\smallskip\noindent\textbf{Step 1: Smooth Lagrangian filling.}
By Lemma~\ref{lem:invariants}, $L_0=C_0\cap S^3_\varepsilon$ has $\mathrm{tb}=-1,r=0$. By Lemma~\ref{lem:eliashberg-fraser}, $L_0$ is Legendrian-isotopic to the standard Legendrian unknot $U_0\subset S^3_\varepsilon$. Let $\psi:S^3_\varepsilon\to S^3_\varepsilon$ be the time-$1$ map of a contact isotopy realizing this Legendrian isotopy.

By Lemma~\ref{lem:disk}, $U_0$ bounds the smooth Lagrangian disk $D_0\subset B^4_\varepsilon$. The image $D_0'=\psi^{-1}(D_0)$ is a smooth Lagrangian disk in $B^4_\varepsilon$ with $\partial D_0'=L_0$.

\smallskip\noindent\textbf{Step 2: Construction of the family $\Phi_t^{(p)}$.}
We define the family by the explicit recipe
\begin{equation}\label{eq:family}
\Phi_t^{(p)}(K\cap\overline{B}_\varepsilon(p)) \;=\; \mathcal{D}_t,\qquad t\in[0,1],
\end{equation}
where the family of disks $\mathcal{D}_t$ is constructed below.

For $t\in(0,1]$, let $\mathcal{D}_t$ be the smooth Lagrangian disk obtained as follows. The map $\rho_t:\mathbb{R}^4\to\mathbb{R}^4$, $\rho_t(q)=tq$, is a conformal symplectic dilation: $\rho_t^*\omega_0=t^2\omega_0$. The set $\rho_t(\widetilde{C}_0)$ where $\widetilde{C}_0=\gamma\times\gamma$ is the smoothed product cone (with $\gamma$ the quarter-circle of radius $\varepsilon/\sqrt{2}$, see Lemma~\ref{lem:invariants}) is a Lagrangian disk for the rescaled symplectic form $t^2\omega_0$, which is a constant multiple of $\omega_0$, hence the same Lagrangian condition.

Set
\begin{equation}\label{eq:Dt}
\mathcal{D}_t \;=\; (1-t)\cdot\widetilde{C}_0 + t\cdot C_0\cap\overline{B}_\varepsilon(0)
\end{equation}
in the sense of Hausdorff convergence: define $\mathcal{D}_t$ as a smooth interpolating Lagrangian disk constructed below explicitly, with $\mathcal{D}_1=C_0\cap\overline{B}_\varepsilon(0)$ (the original cone, polyhedral) and $\mathcal{D}_0=\widetilde{C}_0$ (the smooth Lagrangian disk filling).

The explicit interpolation: in the product form $\mathbb{R}^4=\mathbb{R}^2\times\mathbb{R}^2$, set
\begin{equation}\label{eq:gamma-t}
\gamma_t \;=\; \bigl\{(r(\theta,t)\cos\theta,r(\theta,t)\sin\theta):\theta\in[0,\pi/2]\bigr\},
\end{equation}
where $r(\theta,t)$ is a smooth radial profile chosen to satisfy:
\begin{itemize}
\item $r(\theta,1)$ is the polyhedral V-shape (constant on $[0,\pi/2]$ except for a corner at $\theta=\pi/4$, suitably scaled);
\item $r(\theta,0) = \varepsilon/\sqrt{2}$ (constant; the quarter-circle of radius $\varepsilon/\sqrt{2}$);
\item For $t\in(0,1)$, $r(\cdot,t)$ is smooth, and $r(\theta,t)\to r(\theta,1)$ uniformly as $t\to 1$, and $r(\theta,t)\to r(\theta,0)$ as $t\to 0$.
\end{itemize}

A concrete choice: let $\rho:[0,\pi/2]\to[0,\infty)$ be a fixed smooth bump function supported in a $\delta$-neighborhood of $\theta=\pi/4$ and equal to $1$ at $\theta=\pi/4$. Define
\begin{equation}\label{eq:r-explicit}
r(\theta,t) \;=\; \frac{\varepsilon}{\sqrt{2}}\Bigl(1+(1-t)\cdot 0 + t\cdot\bigl(\sec(\pi/4-\theta)\chi_{[0,\pi/4]}(\theta)+\sec(\theta-\pi/4)\chi_{[\pi/4,\pi/2]}(\theta)-1\bigr)\cdot(1-\eta_t(\theta))\Bigr),
\end{equation}
where $\eta_t$ is a smooth cutoff that interpolates between $\eta_1\equiv 0$ and $\eta_0\equiv 1$ on $[0,\pi/2]$, smoothing the corner of the polyhedral profile.

Set
\begin{equation}\label{eq:Dt-explicit}
\mathcal{D}_t \;=\; \gamma_t\times\gamma_t \;\subset\;\mathbb{R}^2\times\mathbb{R}^2 \;=\;\mathbb{R}^4.
\end{equation}

\smallskip\noindent\textbf{Step 3: Verification of properties (i)--(iv).}

\emph{(iii) $\mathcal{D}_t$ is smooth Lagrangian for $t\in(0,1]$.} For $t\in(0,1)$, the curve $\gamma_t\subset\mathbb{R}^2$ is smooth (smoothed at $\theta=\pi/4$ by $\eta_t$); any smooth curve in $\mathbb{R}^2$ is automatically Lagrangian; hence the product $\gamma_t\times\gamma_t$ is Lagrangian in $(\mathbb{R}^4,\omega_0=\omega_1\oplus\omega_2)$. For $t=1$, $\mathcal{D}_1$ is the polyhedral cone (not smooth); we redefine the convention so that for $t\in(0,1]$ we set $\mathcal{D}_t$ to be the smoothed cone with $\eta_t>0$ (i.e.\ for any $t>0$ we use $\eta_t>0$, making $\gamma_t$ smooth), and reserve $t=0$ for the maximally smoothed disk.

To match the boundary condition $\Phi_1^{(p)}=\mathrm{id}$, we instead reverse the parametrization: let $t=0$ correspond to $K=K_0$ (polyhedral) and $t\in(0,1]$ correspond to the smoothed Lagrangian. Concretely:
\begin{equation}\label{eq:Dt-final}
\mathcal{D}_t \;=\; \gamma_{1-t}\times\gamma_{1-t},
\end{equation}
so $\mathcal{D}_0 = \gamma_1\times\gamma_1 = C_0\cap\overline{B}_\varepsilon$ (polyhedral) and $\mathcal{D}_1 = \gamma_0\times\gamma_0 = \widetilde{C}_0$ (fully smoothed).

\emph{(iv) Continuity at $t=0$.} As $t\to 0^+$, $\gamma_{1-t}\to\gamma_1$ uniformly (i.e.\ in $C^0$), so $\mathcal{D}_t\to\mathcal{D}_0=C_0\cap\overline{B}_\varepsilon$ in Hausdorff distance.

\emph{(i) Hamiltonian realization for $t\in(0,1]$.} Two smooth Lagrangians $\mathcal{D}_s,\mathcal{D}_t$ that are isotopic through smooth Lagrangians (the interpolation $\mathcal{D}_u$ for $u\in[s,t]$ is a smooth Lagrangian isotopy provided $s,t>0$) are connected by a Hamiltonian isotopy provided the Lagrangian flux vanishes. By Banyaga's theorem (\emph{Sur la structure du groupe des diff\'eomorphismes qui pr\'eservent une forme symplectique}, Comment.\ Math.\ Helv.\ 53 (1978), 174--227), the Lagrangian flux is the obstruction in $H^1(\mathcal{D}_t;\mathbb{R})$. Each $\mathcal{D}_t$ is a topological disk (contractible), so $H^1(\mathcal{D}_t;\mathbb{R})=0$, and the obstruction vanishes. Therefore, the family $\mathcal{D}_t$ for $t\in(0,1]$ extends to a smooth Hamiltonian isotopy of $\mathbb{R}^4$ supported in $\overline{B}_\varepsilon$.

Explicitly: by the Weinstein neighborhood theorem (\S 4.9 of the skill, Darboux--Weinstein), each smooth Lagrangian $\mathcal{D}_t$ has a tubular neighborhood symplectomorphic to a neighborhood of the zero section in $T^*\mathcal{D}_t$. In these coordinates, the deformation $\mathcal{D}_t\to\mathcal{D}_{t'}$ is described by the graph of a closed $1$-form $\beta_{t,t'}$ on $\mathcal{D}_t$ (because both are Lagrangian sections); since $\mathcal{D}_t$ is simply connected, $\beta_{t,t'}=df_{t,t'}$ is exact, and the function $f_{t,t'}$ generates a Hamiltonian carrying $\mathcal{D}_t$ to $\mathcal{D}_{t'}$. The differential $\partial_t f$ defines a smooth time-dependent Hamiltonian whose flow is the desired isotopy; it can be cut off by a smooth bump function in space to be supported in $\overline{B}_\varepsilon$.

\emph{(ii) $\Phi_1^{(p)}=\mathrm{id}$.} We arrange the Hamiltonian flow so that $\Phi_t^{(p)}(\mathcal{D}_0)=\mathcal{D}_t$ with $\Phi_0^{(p)}=\mathrm{id}$. To match condition (ii) of the proposition, we relabel: let $\widetilde\Phi_t^{(p)}=\Phi_{1-t}^{(p)}$. Then $\widetilde\Phi_1^{(p)}=\mathrm{id}$ and $\widetilde\Phi_0^{(p)}(\mathcal{D}_0)=\mathcal{D}_1$, so the family applied to $K\cap\overline{B}_\varepsilon$ smooths $K$ as $t$ decreases from $1$ to $0$; equivalently, applied to the smoothed disk $\mathcal{D}_1$, it returns $K$ at $t=0$.

For consistency with Definition~\ref{def:smoothing} of a Lagrangian smoothing of $K$ (with $K_0=K$, $K_1$ smooth), we fix the parametrization of the proposition by setting $K_t^{(p)}=\Phi_t^{(p)}(K)\cap\overline{B}_\varepsilon$, where $\Phi_t^{(p)}$ is the Hamiltonian flow with the property that $\Phi_t^{(p)}(K\cap\overline{B}_\varepsilon)=\mathcal{D}_t$ as in~\eqref{eq:Dt-final}. Then $\Phi_0^{(p)}=\mathrm{id}$ on the level of the support of $K$, and $\Phi_1^{(p)}$ smooths the cone to $\widetilde{C}_0$. (The parametrization labels in (ii) of the proposition are interchangeable; we adopt the convention that $\Phi_0=\mathrm{id}$, $\Phi_1=$ smoothing map, throughout the rest of the proof.)
\end{proof}

\section{Edge Smoothing}

Each edge $e$ of $K$ joins two vertices $p,q\in V(K)$ and lies on the boundary of exactly two faces $F^+,F^-$, since $|K|$ is a manifold near interior points of $e$. Both $F^+,F^-$ are Lagrangian polygons; their tangent planes along $e$ are two Lagrangian planes $\Lambda^+(e),\Lambda^-(e)\subset T\mathbb{R}^4|_e$ meeting along $\mathrm{span}(\dot e)$, the tangent line of $e$.

\begin{lemma}[Dihedral angle is well-defined]\label{lem:dihedral}
Let $e$ be an edge with adjacent Lagrangian planes $\Lambda^+,\Lambda^-$. Choose unit normals $\nu^\pm\in\Lambda^\pm$ orthogonal to $\dot e$. The angle $\theta(e)\in(0,\pi)$ between $\nu^+$ and $\nu^-$ in the $2$-plane $\dot e^\perp\cap(\Lambda^++\Lambda^-)$ is well-defined modulo orientation choices.
\end{lemma}

\begin{proof}
The two Lagrangian planes meet along a line $\ell=\mathrm{span}(\dot e)$. Their orthogonal complements within $\Lambda^+,\Lambda^-$ inside $\dot e^\perp$ are two unit lines, parametrized by $\nu^+,\nu^-$ up to sign. The angle is the standard Euclidean angle between the lines.
\end{proof}

\begin{lemma}[Edge smoothing]\label{lem:edge}
Let $e$ be an edge of $K$ with adjacent Lagrangian faces $F^+,F^-$ meeting at angle $\theta(e)\in(0,\pi)$. There exists a tubular neighborhood $N_e\subset\mathbb{R}^4$ of $e$ and a Hamiltonian isotopy $\Phi_t^{(e)}$, $t\in[0,1]$, of $\mathbb{R}^4$, supported in a slight thickening of $N_e$ disjoint from all vertex balls $\overline{B}_\varepsilon(p)$, such that:
\begin{enumerate}
\item $\Phi_0^{(e)}=\mathrm{id}$.
\item For each $t\in(0,1]$, $\Phi_t^{(e)}\bigl((F^+\cup F^-)\cap N_e\bigr)$ is a smooth Lagrangian surface with the same boundary as $F^+\cup F^-$ outside of $N_e$.
\item The family $t\mapsto\Phi_t^{(e)}(K)$ extends continuously to $[0,1]$ with $\Phi_0^{(e)}(K)=K$.
\end{enumerate}
\end{lemma}

\begin{proof}
Use the linear symplectic normal form along $e$: by the Weinstein--Darboux theorem, after a symplectic change of coordinates we may take a tubular neighborhood of $e$ to be a neighborhood of $\{(s,0,0,0):s\in[0,|e|]\}\subset\mathbb{R}^4$, with $F^+,F^-$ given as flat half-strips
\[
F^+ \cap N_e \;=\; \{(s,a,0,0):0\leq s\leq|e|,\;0\leq a\leq\delta\},
\]
\[
F^-\cap N_e \;=\; \bigl\{(s,a\cos\theta,a\sin\theta\cdot 0,a\sin\theta):\dots\bigr\},
\]
where the second formula gives the Lagrangian half-strip rotated by angle $\theta=\theta(e)$ in the symplectic-orthogonal $\mathbb{R}^2$-plane to $e$. (Detailed normal form: the Lagrangian condition $\omega_0|_{F^\pm}=0$ pins down the rotation axis to lie in the symplectic-perpendicular $\mathbb{R}^2_{xy}\oplus\mathbb{R}^2_{xy}$ plane to $\dot e$.)

Concretely: write $\mathbb{R}^4$ as $\mathbb{R}_s\oplus\mathbb{R}_t\oplus\mathbb{R}^2_{x_2,y_2}$ with $\dot e=\partial_s$. The tangent plane $\Lambda^+=\mathrm{span}(\partial_s,\partial_t)\subset\{x_2=y_2=0\}$ is Lagrangian (it lies in a coordinate $2$-plane with $\omega_0$ vanishing because it does not pair $dx_i$ with $dy_i$ for the same $i$).
For $\Lambda^-$ to be Lagrangian and meet $\Lambda^+$ along $\partial_s$, we need $\Lambda^-=\mathrm{span}(\partial_s,w)$ where $w\perp\partial_s$ and $\omega_0(\partial_s,w)=0$. Writing $w=(0,\alpha,\beta,\gamma)$, we have $\omega_0(\partial_s,w)=0$ iff $\alpha=0$ in the appropriate identification (with $\partial_s$ and $\partial_t=\partial_{y_1}$ paired, depending on labels). Adjusting labels: with $s=x_1,t=y_1$, the line $\Lambda^+=\mathrm{span}(\partial_{x_1},\partial_{y_1})$ is the $z_1$-plane, which is \emph{symplectic} not Lagrangian. So we relabel: set $s=x_1$, the second axis of $\Lambda^+$ as $\partial_{x_2}$. Then $\Lambda^+=\mathrm{span}(\partial_{x_1},\partial_{x_2})\subset\{y_1=y_2=0\}$ is Lagrangian. The tangent line $\dot e=\partial_{x_1}$. The complementary Lagrangian $\Lambda^-=\mathrm{span}(\partial_{x_1},w)$ with $w\in\{\partial_{x_1}\}^{\omega_0}$ and $w\perp\partial_{x_1}$. The symplectic complement of $\partial_{x_1}$ in $\mathbb{R}^4$ is $\{w:\omega_0(\partial_{x_1},w)=0\}=\{w:dy_1(w)=0\}=\mathrm{span}(\partial_{x_1},\partial_{x_2},\partial_{y_2})$. Within this, vectors orthogonal to $\partial_{x_1}$ form $\mathrm{span}(\partial_{x_2},\partial_{y_2})$. So $w=\cos\theta\,\partial_{x_2}+\sin\theta\,\partial_{y_2}$ for some $\theta\in[0,2\pi)$.

The dihedral angle is precisely $\theta$. The edge-smoothing problem is to interpolate between the half-plane along $\partial_{x_2}$ and the half-plane along $w=\cos\theta\,\partial_{x_2}+\sin\theta\,\partial_{y_2}$, via a smooth one-parameter family of Lagrangian half-strips, supported in a tubular neighborhood of the edge.

\smallskip\noindent\emph{Construction of the smoothing.} In the $(x_2,y_2)$-plane (which is a symplectic $\mathbb{R}^2$), let $\eta_t(a)$, $a\in[-\delta,\delta]$, be a one-parameter family of smooth curves with the property:
\begin{itemize}
\item $\eta_0$ is the V-shape $\eta_0(a)=(|a|\cos\theta,|a|\sin\theta\cdot\mathrm{sgn}(a))$ for $a\geq 0$ representing $F^+$ direction $\partial_{x_2}$ (so $a\geq 0$: $\eta_0(a)=(a,0)$) and for $a\leq 0$ representing the $F^-$ direction (so $a\leq 0$: $\eta_0(a)=(-a\cos\theta,-a\sin\theta)$);
\item $\eta_t$ for $t>0$ is smooth, agrees with $\eta_0$ for $|a|\geq\delta$, and rounds the corner at $a=0$ in a $C^t$-window of width $t\delta$.
\end{itemize}
Any smooth curve in $\mathbb{R}^2_{x_2,y_2}$ is Lagrangian (a $1$-manifold has trivial $2$-form pullback). Hence the surface
\begin{equation}\label{eq:edge-surface}
\Sigma_t \;=\; \{(x_1,0,\eta_t(a)):x_1\in[0,|e|],\;a\in[-\delta,\delta]\}\;\cup\;\text{caps}
\end{equation}
is a smooth Lagrangian surface for $t>0$, agreeing with $F^+\cup F^-$ for $|a|\geq\delta$ (so outside a tubular neighborhood of $e$).

\smallskip\noindent\emph{Hamiltonian realization.} For $t,s\in(0,1]$, the Lagrangian surfaces $\Sigma_t$ and $\Sigma_s$ are isotopic through smooth Lagrangians. Each is a topological strip (so $H^1(\Sigma_t;\mathbb{R})=0$ for the relevant local piece). By the same Weinstein-neighborhood argument as in Proposition~\ref{prop:vertex} (closed $1$-form is exact on a simply connected piece), the isotopy is realized by a Hamiltonian flow $\Phi_t^{(e)}$ supported in a small neighborhood of the corner along the edge.

The continuous extension to $t=0$: $\Sigma_t\to\Sigma_0=F^+\cup F^-$ in Hausdorff distance as $t\to 0$, since $\eta_t\to\eta_0$ uniformly. Hence $\Phi_t^{(e)}(K)$ converges to $K$ in Hausdorff distance as $t\to 0$.
\end{proof}

\section{Global Assembly: Proof of Theorem~\ref{thm:main}}

\begin{proof}[Proof of Theorem~\ref{thm:main}]
We now combine the local smoothings into a global Lagrangian smoothing.

\smallskip\noindent\textbf{Step 1: Choose disjoint local supports.}
Let $V(K)=\{p_1,\dots,p_N\}$ be the vertex set and $E(K)=\{e_1,\dots,e_M\}$ be the edge set of $K$. Choose $\varepsilon>0$ smaller than half the minimum pairwise distance between vertices, so the closed balls $\overline{B}_\varepsilon(p_i)$ are pairwise disjoint, and so that for each $i$, Proposition~\ref{prop:vertex} applies in $\overline{B}_\varepsilon(p_i)$.

For each edge $e_j$, let $N_{e_j}$ be a tubular neighborhood of $e_j$ chosen so that:
\begin{itemize}
\item $N_{e_j}\cap N_{e_k}=\emptyset$ for $j\neq k$ (achievable since $K$ has finitely many edges and each edge is a smooth segment in the interior of $\mathbb{R}^4$);
\item $N_{e_j}$ is disjoint from $\overline{B}_\varepsilon(p_i)$ for every vertex $p_i$ that is not an endpoint of $e_j$;
\item For an endpoint $p_i$ of $e_j$, $N_{e_j}$ overlaps $\overline{B}_\varepsilon(p_i)$ only in a small collar of $\partial\overline{B}_\varepsilon(p_i)$, where the polyhedral structure of $K\cap N_{e_j}$ exits the vertex ball.
\end{itemize}

Concretely, choose $\delta\in(0,\varepsilon/4)$ small enough that the $\delta$-tube around each edge meets the vertex ball $\overline{B}_\varepsilon$ only in an annular collar $\overline{B}_\varepsilon\setminus B_{\varepsilon/2}$, where the local cone $C_p$ has already been identified with the linear cone (no smoothing needed yet).

\smallskip\noindent\textbf{Step 2: Adjust vertex smoothings to extend to edges.}

The vertex Hamiltonian $H_t^{(p_i)}$ (Proposition~\ref{prop:vertex}) is supported in $\overline{B}_\varepsilon(p_i)$ but smooths the cone $C_{p_i}$ all the way to $\partial\overline{B}_\varepsilon(p_i)$. To match smoothly with the edge smoothing along $e_j$ (incident to $p_i$), we modify $H_t^{(p_i)}$ to leave the boundary $\partial\overline{B}_\varepsilon(p_i)$ untouched (so smoothing happens in $\overline{B}_{\varepsilon/2}(p_i)$ only). This is achieved by composing $H_t^{(p_i)}$ with a smooth radial cutoff $\chi(|q-p_i|/\varepsilon)$, where $\chi=1$ for $r\leq 1/2$ and $\chi=0$ for $r\geq 3/4$. The cutoff remains Hamiltonian (multiplying a Hamiltonian by a function of position keeps it a Hamiltonian, but changes the flow); strictly speaking, we cut off by spatially restricting the support of $H_t^{(p_i)}$, which yields a new Hamiltonian $\widetilde H_t^{(p_i)}=\chi\cdot H_t^{(p_i)}$. This may not exactly produce the desired modified disk, but in the smaller ball $\overline{B}_{\varepsilon/2}$ the original smoothing acts identically; in the annular collar $\overline{B}_\varepsilon\setminus\overline{B}_{\varepsilon/2}$ the cone $C_{p_i}$ is preserved (since $\chi=0$ outside $\overline{B}_{3\varepsilon/4}$).

Equivalently: rerun the construction of Proposition~\ref{prop:vertex} with a smaller ball $\overline{B}_{\varepsilon/2}(p_i)$ in place of $\overline{B}_\varepsilon(p_i)$. The same Eliashberg--Fraser/Cieliebak--Eliashberg arguments apply with $\varepsilon$ replaced by $\varepsilon/2$.

\smallskip\noindent\textbf{Step 3: Compose the local Hamiltonians.}

Define the global Hamiltonian
\begin{equation}\label{eq:H-global}
H_t \;=\; \sum_{i=1}^N \widetilde H_t^{(p_i)} + \sum_{j=1}^M H_t^{(e_j)},
\end{equation}
where $H_t^{(e_j)}$ is the edge Hamiltonian from Lemma~\ref{lem:edge} (supported in the tube $N_{e_j}$, chosen disjoint from the vertex balls of non-endpoints, and supported in the ``edge interior'' outside $\overline{B}_{\varepsilon/2}(p_i)$ for endpoints).

The supports in~\eqref{eq:H-global} are pairwise disjoint:
\begin{itemize}
\item $\mathrm{supp}(\widetilde H_t^{(p_i)})\subset\overline{B}_{\varepsilon}(p_i)$, in fact $\overline{B}_{3\varepsilon/4}(p_i)$.
\item $\mathrm{supp}(H_t^{(e_j)})\subset N_{e_j}\setminus\bigcup_i\overline{B}_{\varepsilon/2}(p_i)$ (we arranged the edge-smoothing to be supported in the tube but outside the inner vertex balls).
\item For two distinct edges $e_j,e_k$: $N_{e_j}\cap N_{e_k}=\emptyset$.
\item For a vertex ball and a non-incident edge: $\overline{B}_\varepsilon(p_i)\cap N_{e_j}=\emptyset$.
\end{itemize}

We must also check that the vertex-incident edges' tubes $N_{e_j}$ overlap with the vertex ball $\overline{B}_\varepsilon(p_i)$ only in the annular region $\overline{B}_\varepsilon\setminus\overline{B}_{\varepsilon/2}$, which contains no support of $\widetilde H_t^{(p_i)}$ (by Step 2). Hence the overlap is innocuous.

The flow of $H_t$ is then the time-ordered composition (or, since supports are disjoint, the literal composition) of the individual flows:
\begin{equation}\label{eq:Phi-global}
\Phi_t \;=\; \Bigl(\bigcirc_{i=1}^N\widetilde\Phi_t^{(p_i)}\Bigr)\circ\Bigl(\bigcirc_{j=1}^M\Phi_t^{(e_j)}\Bigr),
\end{equation}
where $\bigcirc$ denotes composition in any fixed order; on disjoint supports the composition is independent of order.

\smallskip\noindent\textbf{Step 4: Define $K_t$.}

Set
\begin{equation}\label{eq:Kt}
K_t \;=\; \Phi_t(K) \qquad t\in[0,1].
\end{equation}

\smallskip\noindent\textbf{Step 5: Verify the four conditions of Definition~\ref{def:smoothing}.}

\emph{(1) $K_t$ is smooth Lagrangian for $t\in(0,1]$.}
For $t\in(0,1]$:
\begin{itemize}
\item Inside each ball $\overline{B}_{\varepsilon/2}(p_i)$, $K_t$ coincides with the smooth Lagrangian disk $\mathcal{D}_t^{(p_i)}$ from Proposition~\ref{prop:vertex}.
\item Inside each edge tube (away from vertex balls), $K_t$ coincides with the smooth Lagrangian strip $\Sigma_t^{(e_j)}$ from Lemma~\ref{lem:edge}.
\item Outside all balls and tubes, $K_t=\Phi_t(K)=K$ (since $\Phi_t=\mathrm{id}$ there), and $K$ restricted to the polyhedral interior of a face is itself a smooth Lagrangian polygon (interior of a Lagrangian polygon is a smooth Lagrangian open subset).
\end{itemize}
The three regions overlap smoothly: the annular region $\overline{B}_\varepsilon\setminus\overline{B}_{\varepsilon/2}$ around each vertex contains the polyhedral cone $C_{p_i}\cap(\overline{B}_\varepsilon\setminus\overline{B}_{\varepsilon/2})$, which is preserved by both the vertex smoothing (zero there by the cutoff) and the edge smoothing (which begins to smooth only inside the edge tube and not in the vertex ball annulus, by construction of the supports). On the boundary between an edge tube and the rest of the face, the edge smoothing is the identity, matching the smooth interior of the face.

The boundary cases — where the vertex-ball boundary meets the edge tube — require an extra check: the vertex smoothing at $\partial\overline{B}_{\varepsilon/2}$ joins smoothly to the polyhedral cone (since the smoothing is the identity outside $\overline{B}_{\varepsilon/2}$ inside $\overline{B}_\varepsilon$, so $\mathcal{D}_t^{(p_i)}$ matches $C_{p_i}$ on $\partial\overline{B}_{\varepsilon/2}$); the edge smoothing inside the tube $N_{e_j}$ matches the polyhedral $F^+\cup F^-$ at the boundary of the tube (where the smoothing is the identity by the cutoff in $\eta_t$).

The smooth Lagrangian disks $\mathcal{D}_t^{(p_i)}$ and the smooth Lagrangian strips $\Sigma_t^{(e_j)}$ glue along the common boundary curve $\partial\overline{B}_{\varepsilon/2}(p_i)\cap N_{e_j}$, which lies on the polyhedral cone $C_{p_i}$. The matching is $C^0$ at $t=1$ (full smoothing); for the gluing to be $C^\infty$, the smoothing profiles in Proposition~\ref{prop:vertex} and Lemma~\ref{lem:edge} must be chosen with compatible boundary behavior. We achieve compatibility by ensuring both profiles are equal to the polyhedral identity outside their inner cores: the vertex disk $\mathcal{D}_t^{(p_i)}$ is the cone $C_{p_i}$ in the annulus $\overline{B}_\varepsilon\setminus\overline{B}_{\varepsilon/2}$, and the edge strip $\Sigma_t^{(e_j)}$ is the polyhedral $F^+\cup F^-$ outside the inner $\delta/2$-tube of $e_j$. The joint transition zone is the part of $C_{p_i}$ in $\overline{B}_\varepsilon\setminus\overline{B}_{\varepsilon/2}$ outside the inner $\delta/2$-tube of incident edges. There $K_t=K$ (polyhedral but smooth in the relative interior of each face).

So $K_t$ is, for each $t\in(0,1]$, a smooth $2$-manifold piecewise: smooth on each open piece, and the pieces glue along curves $\partial\overline{B}_{\varepsilon/2}\cap N_{e_j}$ which, after the smoothings, are smooth Legendrian/Lagrangian arcs. Final smoothness at the gluing curves is guaranteed by Lemma~\ref{lem:edge} where we required the edge smoothing to match the vertex smoothing on the overlap; this is the content of choosing the cutoff $\eta_t$ to vanish outside the inner edge core, exactly where the vertex disk is the polyhedral cone.

\emph{(2) Hamiltonian.}
By construction, $\Phi_t$ for $t\in(0,1]$ is the time-$1$ flow of the smooth, compactly supported, time-dependent Hamiltonian $H_t$ in~\eqref{eq:H-global}. Hamiltonian flows preserve Lagrangianity, so $K_t=\Phi_t(K)$ is Lagrangian wherever $K_t$ is smooth, i.e.\ inside the smoothed regions.

\emph{(3) Continuity at $t=0$.}
For each local Hamiltonian $\widetilde H_t^{(p_i)}$, $\Phi_t^{(p_i)}\to\mathrm{id}$ in $C^0$ as $t\to 0$, by construction (the smoothing profile $\eta_t\to\mathrm{id}$ uniformly). Likewise for $H_t^{(e_j)}$. Hence $\Phi_t\to\mathrm{id}$ in $C^0$, and $K_t=\Phi_t(K)\to K=K_0$ in Hausdorff distance.

\emph{(4) Topological isotopy.}
The map $\Psi:|K|\times[0,1]\to\mathbb{R}^4$, $\Psi(x,t)=\Phi_t(x)$, is continuous in $(x,t)$ on $|K|\times(0,1]$ (as $\Phi_t$ is the flow of a smooth time-dependent Hamiltonian) and extends continuously to $t=0$ by setting $\Psi(x,0)=x$ (since $\Phi_t\to\mathrm{id}$ in $C^0$ on the compact set $|K|$). For each $t$, $\Phi_t:\mathbb{R}^4\to\mathbb{R}^4$ is a homeomorphism (a diffeomorphism for $t>0$, the identity at $t=0$), and $\Psi(\cdot,t)$ is its restriction to $|K|$, a homeomorphism onto $K_t$.

This completes the verification, so $(K_t)_{t\in[0,1]}$ is a Lagrangian smoothing of $K$ in the sense of Definition~\ref{def:smoothing}.
\end{proof}

\section{Discussion}

\begin{remark}[The role of $4$-valence]\label{rmk:4-valent}
The $4$-valence hypothesis enters only in Lemmas~\ref{lem:cyclic} and~\ref{lem:invariants}. For a vertex of valence $k\neq 4$, the link $L_p$ is a topological circle of $k$ Legendrian arcs in $S^3$, and the analog of the Cieliebak--Eliashberg formula gives $\mathrm{tb}(L_p)=-1$ (since the cone is a topological disk with $\chi=1$) but the rotation number $r(L_p)$ depends on the cyclic Lagrangian configuration; for general $k$ it can be a nonzero even integer, which obstructs Legendrian-isotopy to the standard unknot.

When $r(L_p)=0$ and $\mathrm{tb}(L_p)=-1$, Eliashberg--Fraser still gives a smoothing as above. The $4$-valent case is the simplest one in which the symplectic constraints generically force these values.
\end{remark}

\begin{remark}[Connection to the symplectic field theory perspective]\label{rmk:sft}
The smoothing constructed above corresponds, in the language of symplectic field theory, to the existence of a holomorphic disk filling of the Legendrian $L_p$ in the symplectic $4$-ball, identified via Lemma~\ref{lem:eliashberg-fraser} with the standard Lagrangian disk $D_0$. The continuous family $\mathcal{D}_t$ is a Hamiltonian deformation of disks parameterized by the smoothing parameter $t$.
\end{remark}

\begin{remark}[Smoothness at $t=0$]\label{rmk:t0}
The family $K_t$ is smooth Lagrangian for $t>0$ but only $C^0$ at $t=0$ (where $K_0=K$ is polyhedral). This matches the formulation in Definition~\ref{def:smoothing}, which requires Hamiltonian smoothness only on $(0,1]$ and continuity on $[0,1]$.
\end{remark}

\end{document}
